\hypertarget{cont3:paper-ix-compatible-polynomial-towers-and-p-adic-roots}{%
\subsection{Paper IX: compatible polynomial towers and p-adic
roots}\label{cont3:paper-ix-compatible-polynomial-towers-and-p-adic-roots}}

\textbf{Author: Komorebi P. --- 8 September 2026.}

This block completes the nonsingular branch of the passage from one-step
polynomial lifting to an infinite compatible tower and an actual p-adic
root. Surjectivity is required only of the reduced Jacobian at the
initial root. The proof derives its persistence, constructs the tower,
proves convergence and verifies the original polynomial equations at the
limit.

\hypertarget{cont3:main-theorem}{%
\subsubsection{Main theorem}\label{cont3:main-theorem}}

Let \(p\) be prime, \(L\ge1\), and \[
F\in\mathbb Z[X_1,\ldots,X_d]^e,\qquad x\in\mathbb Z^d,
\qquad F(x)\equiv0\pmod{p^L}.
\] Assume the actual reduced Jacobian \[
A_x=DF(x)\bmod p:\mathbb F_p^d\longrightarrow\mathbb F_p^e
\] is surjective. Then Lean proves that there is a point \[
\boxed{a\in\mathbb Z_p^d,\qquad F(a)=0,\qquad
 a\equiv x\pmod{p^L}.}
\] Here the coefficients of the original integer polynomials are cast
into \(\mathbb Z_p\), and their multivariable evaluation is exactly
zero. This is an existence theorem for a fixed polynomial presentation
with the stated Jacobian hypothesis.

The stronger checked statement constructs \(z_n\in\mathbb Z^d\),
\(n\ge0\), with \[
\begin{aligned}
z_0&=x,\\
F(z_n)&\equiv0\pmod{p^{L+n}},\\
z_{n+1}&\equiv z_n\pmod{p^{L+n}},\\
DF(z_n)\bmod p&=A_x,\\
z_n&\longrightarrow a\quad\text{in }\mathbb Z_p^d,\\
a-z_n&\in p^{L+n}\mathbb Z_p^d.
\end{aligned}
\] The last inclusion retains every prescribed precision. It is a
divisibility lower bound, not an assertion that the error has exactly
that valuation. An approximation can already equal the limiting root.

\hypertarget{cont3:why-the-single-initial-hypothesis-suffices}{%
\subsubsection{Why the single initial hypothesis
suffices}\label{cont3:why-the-single-initial-hypothesis-suffices}}

The one-step Taylor theorem already supplies the exact correction
equation \[
b_n+A_{z_n}t=0,\qquad
b_n=F(z_n)/p^{L+n}\bmod p.
\] Surjectivity supplies a correcting digit for every residual \(b_n\).
The chosen correction gives an actual new integer point \(z_{n+1}\),
using the proved nonlinear Taylor remainder.

Every coordinate change is divisible by \(p^{L+n}\), hence by \(p\).
Each formal partial derivative is itself an integer polynomial, so its
value is unchanged modulo \(p\). Therefore \[
A_{z_{n+1}}=A_{z_n}=A_x.
\] The proof assumes no separate surjectivity or solvability statement
at later levels. A dependent recursion chooses each next point from the
proved one-step existence theorem. It produces one compatible sequence,
rather than unrelated witnesses at different precisions.

The integer tower construction works more generally for every positive
modulus \(q\), including composite \(q\) and \(q=1\), with a surjective
map over \(\mathbb Z/q\mathbb Z\). The p-adic completion and the
rank-based branching formulas in this block use prime \(p\).

\hypertarget{cont3:convergence-and-exact-polynomial-vanishing}{%
\subsubsection{Convergence and exact polynomial
vanishing}\label{cont3:convergence-and-exact-polynomial-vanishing}}

Compatibility implies, for \(m\ge n\), \[
p^{L+n}\mid z_m-z_n.
\] After casting to \(\mathbb Z_p\), this bounds each coordinate
distance by \(p^{-(L+n)}\), so the sequence is Cauchy. Completeness
supplies its limit. Closed norm bounds show that the limit retains the
congruence at every stage.

The generic convergence theorem is proved independently of the
polynomial application. It also proves \[
p^{L+n}\mid v_n\ \text{for all }n
\quad\Longrightarrow\quad v_n\longrightarrow0\text{ in }\mathbb Z_p.
\] Apply this to each integer sequence \(F_i(z_n)\). Continuity of
multivariable polynomial evaluation also gives \(F_i(z_n)\to F_i(a)\).
Uniqueness of limits then yields \(F_i(a)=0\).

The cast-evaluation identity is explicitly checked, so the continuous
function is the evaluation of the original polynomials. Convergence and
polynomial vanishing are conclusions, not assumptions of the main
existence theorem.

\hypertarget{cont3:exact-number-of-choices-at-each-level}{%
\subsubsection{Exact number of choices at each
level}\label{cont3:exact-number-of-choices-at-each-level}}

Surjectivity implies \[
\operatorname{rank}_{\mathbb F_p}A_x=e,\qquad e\le d.
\] Combining this with the preceding exact residue-lift count gives \[
\boxed{\#\mathcal L_{z_n,L+n}=p^{d-e}\quad\text{for every }n.}
\] These are distinct solutions modulo \(p^{L+n+1}\) lying above the
selected point modulo \(p^{L+n}\). The original finite obstruction group
at each smooth point has cardinality one.

The theorem constructs one path through these finite choices. When
\(d>e\), it does not claim that all paths have the same limit. The
generic uniqueness theorem concerns the limit of a fixed sequence only.

\hypertarget{cont3:boundary-controls}{%
\subsubsection{Boundary controls}\label{cont3:boundary-controls}}

The positive control theorems and rejected invalid proofs cover three
distinctions.

\begin{enumerate}
\def\labelenumi{\arabic{enumi}.}
\tightlist
\item
  \textbf{An initial lift does not guarantee an infinite tower without
  the Jacobian condition.} \(X^2-5\) has the odd lower root modulo two
  and a lift modulo four, but no root modulo eight. Lean proves that no
  integer tower can continue through all precisions in this example.
\item
  \textbf{Surjective rectangular Jacobians do not give uniqueness of all
  roots.} The system \(F(X,Y)=X\) has a surjective Jacobian and distinct
  dyadic roots \((0,0)\) and \((0,2)\), both in the same residue modulo
  two.
\item
  \textbf{The retained initial precision cannot be increased without
  justification.} For \(F(X)=X-2\), zero is a lower root modulo two with
  surjective derivative, but its exact root two does not retain zero
  modulo four.
\end{enumerate}

The controls are theorem checks of these explicit examples. They
supplement the universal proof audit and do not replace it.

For several polynomial blocks, the hypothesis applies to their combined
system. Surjectivity of each block separately does not imply
surjectivity of the combined Jacobian or shared liftability. The
preceding shared-correction obstruction remains relevant in that
situation.
